% vim: ai sw=2 spell spelllang=cs

\begin{frame}
  \frametitle{Komunikace s HW}

  \heading{LDD3 kap. 9 a 12}

  \iffimuni
    \bigskip

    \tableofcontents
  \fi

\end{frame}

\mysection{I/O}

\begin{frame}
  \frametitle{I/O}
  \heading{Na x86: 2 přístupy}
  \begin{itemize}
  \item Porty
  \item Memory Mapped I/O
  \end{itemize}

\end{frame}

\iffalse
\subsection{I/O -- porty}
\fi

\begin{frame}{I/O}{Porty}

  \heading{Porty}
  \begin{itemize}
    \item Speciální adresový prostor
    \begin{itemize}
      \item Závislé na architektuře
      \item Speciální instrukce (\code{in}, \code{out} na x86)
    \end{itemize}

    \pause

    \item Samostatná (malá) sběrnice
    \begin{itemize}
      \item Na x86: řadič klávesnice, PC spkr, staré časovače a ovladače
      přerušení, ladicí port (\code{0x80} $\Rightarrow$ segmentový displej
      na desce), \dots
    \end{itemize}
  \end{itemize}

  \pause

  \heading{API}
  \begin{itemize}
  \item \header{linux/io.h}, \header{linux/ioport.h}
  \item Požadavek: \code{request_region}, \code{release_region}
  \item R/W: \code{inX(port)}, \code{outX(co, port)}, kde
    $\text{\code{X}}\in\left\{\text{\code{b}},\text{\code{w}},\text{\code{l}}\right\}$
  \end{itemize}

  \medskip
  \pause

  \heading{Demo:} pb173/08 \\
  \heading{Úkol:} Přečíst port 0x80, zapsat do něj 1B a znovu přečíst

\end{frame}

\iffalse
\subsection{I/O -- MMIO}
\fi

\begin{frame}{I/O}{MMIO}

  \heading{Memory Mapped I/O}
  \begin{itemize}
    \item Součástí fyzického adresového prostoru
    \item Přístup standardním čtením/zápisem
    \begin{itemize}
      \item Nutnost přemapovat na virtuální adresy
    \end{itemize}
  \end{itemize}

  \pause

  \heading{API}
  \begin{itemize}
    \item \header{linux/io.h}, \header{linux/ioport.h}

    \item Požadavek: \code{request_mem_region}, \code{release_mem_region}
    \begin{itemize}
      \item \file{/proc/iomem}
    \end{itemize}

    \pause

    \item Mapování: \code{virt = ioremap(phys)}, \code{iounmap(virt)}
    \begin{itemize}
      \item \code{/proc/vmallocinfo} (ne stará jádra)
    \end{itemize}

    \pause

    \item R/W: \code{readX(odkud)}, \code{writeX(co, kam)}, kde
      $\text{\code{X}}\in\left\{\text{\code{b}},\text{\code{w}},\text{\code{l}},\text{\code{q}}\right\}$
    \begin{itemize}
      \item \pozor: nelze číst jako z paměti (optimalizace, pořadí bytů, \dots)
    \end{itemize}
  \end{itemize}

\end{frame}

\mysection{PCI zařízení}

\begin{frame}
  \frametitle{PCI zařízení}

  \begin{itemize}
    \item PCI, PCI-X, PCIe

    \pause

    \item Hierarchická sběrnice
    \begin{itemize}
      \item Identifikace doména:bus:slot.funkce
      \item Bridge (=routery)
    \end{itemize}

    \pause

    \item Konfigurační prostor
    \begin{itemize}
      \item Automatická konfigurace

      \item ID zařízení (vendor, device), I/O prostory, IRQ
      \begin{itemize}
	\item \emph{Obsah} I/O prostorů -- specifikace zařízení (výrobce)
      \end{itemize}

      \item \cmd{lspci}
    \end{itemize}

    \pause

    \item Podrobnosti v PCI specifikaci
  \end{itemize}

\end{frame}

\begin{frame}
  \frametitle{PCI zařízení v jádře}

  \begin{itemize}
    \item \header{linux/pci.h}, \doc{Documentation/pci/*}
    \item \code{struct pci_dev}, \code{struct pci_bus}

    \pause

    \item \code{pci_\{set,get\}_drvdata}
    \begin{itemize}
      \item Uloží/načte programátorova data do/z \code{struct pci_dev}
    \end{itemize}

    \pause

    \item Referenční počet
    \begin{itemize}
      \item PCI zařízení může být (fyzicky) odebráno ze systému
      \item Ale \code{struct pci_dev} zůstává, pokud reference $> 0$

      \pause

      \item Zvýšení: \code{pci_dev_get}
      \item Snížení: \code{pci_dev_put}
    \end{itemize}

    \pause

    \item Identifikace (ID) zařízení
    \begin{itemize}
      \item \code{pci_dev->vendor}, \code{pci_dev->device}, \dots
      \item \code{PCI_ANY_ID} značí jakékoliv ID (při hledání)
    \end{itemize}
  \end{itemize}

\end{frame}

\begin{frame}[fragile]
  \frametitle{Hledání zařízení}
  \framesubtitle{Iterátory (starší)}

  \begin{itemize}
    \item Dovoluje přístup k zařízení více ovladači
    \item Nepodporuje hotplug
    \item \code{pci_get_device} (vendor, device)
  \end{itemize}

  \pause

  \begin{block}{Iterace přes zařízení}
    \begin{lstlisting}[aboveskip=2pt,belowskip=2pt]
struct pci_dev *pdev = NULL;

while ((pdev = pci_get_device(VENDOR, DEVICE, pdev))) {
  pr_info("%.2x:%.2x.%.2x\n",
          pdev->bus->number,
          PCI_SLOT(pdev->devfn),
          PCI_FUNC(pdev->devfn));
}
    \end{lstlisting} %stopzone
  \end{block}

  \pause
  \smallskip

  \heading{Úkol:} vypsat všechna zařízení v systému (jejich
  \emph{vendor+device} ID)

\end{frame}

\begin{frame}[fragile]
  \frametitle{Hledání zařízení}
  \framesubtitle{Událostmi}

  \begin{itemize}
    \item Registrace háčků a seznamu chtěných zařízení
    \begin{itemize}
      \item Seznam: \code{struct pci_device_id} (vendor, device, atd.)
      \item Háčky: \code{struct pci_driver} (probe, remove, suspend, atd.)
    \end{itemize}
    \item \code{pci_register_driver}, \code{pci_unregister_driver}
  \end{itemize}

  \begin{block}{\footnotesize Události a PCI zařízení}
    \begin{lstlisting}[aboveskip=2pt,belowskip=2pt]
static struct pci_device_id my_table[] = {
  { PCI_DEVICE(VENDOR1, DEVICE1) }, { PCI_DEVICE(VENDOR1, DEVICE2) },
  { PCI_DEVICE(0x8086, PCI_ANY_ID), .driver_data = 1 },       { 0, }
};
MODULE_DEVICE_TABLE(pci, my_table);
static struct pci_driver my_pci_driver = {
  .name = "my_driver",   .id_table = my_table,   .probe = my_probe,
};
int my_probe(struct pci_dev *pdev, const struct pci_device_id *id) {
  pr_info("%2.x %lu\n", pdev->bus->number, id->driver_data);
  return 0;
}
    \end{lstlisting} %stopzone
  \end{block}

\end{frame}

\begin{frame}{Úkol}

  \heading{Navázání PCI zařízení}
  \begin{enumerate}
    \item \header{linux/pci.h}

    \item Definujte \code{struct pci_device_id}
    \begin{itemize}
      \item \cmd{lspci -nn} a najděte EDU a jeho vendor+device ID
    \end{itemize}

    \item Definujte \code{struct pci_driver}
    \begin{itemize}
      \item \code{probe}, \code{remove}, \code{name}, \code{id_table}
    \end{itemize}

    \item Definujte háčky (viz definici \code{struct pci_driver})
    \begin{itemize}
      \item \scriptsize{\code{int (*probe)(struct pci_dev *, const struct pci_device_id *)}}
      \item \scriptsize{\code{void (*remove)(struct pci_dev *)}}
    \end{itemize}

    \item V \code{probe} a \code{remove} vypište
    \begin{itemize}
      \item \code{pdev->bus->number}
      \item \code{PCI_SLOT(pdev->devfn)}
      \item \code{PCI_FUNC(pdev->devfn)}
    \end{itemize}

    \item Zavolejte \code{pci_register_driver} a \code{pci_unregister_driver}
    \item Vložte a odeberte modul a zkontrolujte \cmd{dmesg}
  \end{enumerate}

\end{frame}

\begin{frame}{Práce s PCI}

  \heading{Probe}
  \begin{itemize}
    \item Inicializace PCI zařízení
    \begin{itemize}
      \item \code{pci_enable_device}
      \item Do té doby nelze některé vlastnosti \code{struct pdev} používat
        (irq)
    \end{itemize}

    \pause

    \item Rezervace a mapování I/O (jako v první části \iffimuni cvičení\else
      sekce\fi)
    \begin{itemize}
      \item Požadavek: \code{pci_request_region}, \code{pci_request_regions}

      \pause

      \item Mapování: \code{pci_ioremap_bar} -- alias pro
	\code{ioremap(pci_resource_start(), pci_resource_len())}
    \end{itemize}

    \pause

    \item Nastavení/detekce zařízení
      \begin{itemize}
	\item Závisí na zařízení
      \end{itemize}
  \end{itemize}

  \pause
  \medskip

  \heading{Remove}
  \begin{itemize}
    \item Opak Probe
    \item Tj.\ \code{iounmap}, \code{pci_release_region}(\code{s}),
      \code{pci_disable_device}
  \end{itemize}

\end{frame}

\mysection{Zobecnění na jiné sběrnice}

\begin{frame}
  \frametitle{Zobecnění}

  \heading{Většina ostatních sběrnic funguje stejně}
  \begin{itemize}
    \item Např.\ USB, I$^2$C, HID, IEEE1394, ACPI, INPUT, EISA, \dots
    \item Nějaké probe/remove
    \item Seznam ID
    \item Registrace ,,ovladače''
  \end{itemize}

\end{frame}

\mysection{EDU karta}

\begin{frame}
  \frametitle{Specifikace baru 0 karty EDU}

\begin{block}{Specifikace baru 0 karty EDU}
%\footnotesize
\begin{center}
\begin{tabular}{|l|l|l|l|} \hline
Offset & Len & R/W & Meaning \\ \hline\hline
0x0000 & 4B  & R   & ID \& Revision \\ \hline
0x0004 & 4B  & R/W & Inverted value \\ \hline
0x0008 & 4B  & R/W & Factorial \\ \hline
0x0020 & 4B  & R/W & Status    \\ \hline
\end{tabular}
\end{center}
\end{block}

  \begin{itemize}
    \item ID \& Revision: \code{0xRRrr00edu}, RR -- major, rr -- minor
    \item Inverted value: zapsané číslo se obrátí (C operátor \code{\~})
    \item Faktoriál: vypočte se faktoriál (je nutné počkat na status bit 0)
    \item Status:
    \begin{itemize}
      \item Bit 0 -- probíhá výpočet faktoriálu
      \item Bit 7 -- má faktoriál generovat přerušení?
    \end{itemize}
  \end{itemize}

\end{frame}

\begin{frame}
  \frametitle{Úkol}

  \heading{Vyčtení identifikace karty EDU a práce s ní}
  \begin{enumerate}
    \item Povolte zařízení (\code{pci_enable_device})
    \item Rezervujte bar 0 (\code{pci_request_region})
    \item Vypište fyzickou adresu baru 0 a porovnejte s výstupem \cmd{lspci}
    \item Přemapujte bar 0 (\code{ioremap})

    \item Přečtěte a \emph{rozkódujte} identifikaci a revizi (\code{readX})
    \begin{itemize}
      \item Vypište zvlášť major a minor revizi
    \end{itemize}

    \item Ověřte živost karty pomocí invertovacího registru

    \iffimuni\else
    \item Vypočtěte kartou faktoriál 3, 4 a 5 a vypište
    \fi

    \item Ukliďte v remove (unmap, release, disable)
  \end{enumerate}

\end{frame}
